\chapter{\IfLanguageName{dutch}{Proof-of-Concept}{Proof-of-Concept}}%
\label{ch:poc}

\section{Architectuur van de applicatie}

De proof-of-concept werd geïmplementeerd als een webapplicatie die HLASM-broncode omzet naar een control-flowdiagram in Mermaid.js-syntax.
Omdat het een webapplicatie is, hoeft de gebruiker niets te installeren.  
De tool kan meteen gebruikt worden via de browser. Dit maakt het eenvoudig om de applicatie snel te testen en te gebruiken.

De werking van de applicatie bestaat uit drie stappen:

\begin{enumerate}
    \item \textbf{Parsing}: de ingevoerde HLASM-broncode wordt opgesplitst in individuele instructies. Per regel worden label, mnemonic en operanden geïdentificeerd op basis van de kolomgebaseerde structuur.
    
    \item \textbf{Control-flow analyse}: op basis van de gedetecteerde instructies wordt bepaald hoe de uitvoering van het programma kan verlopen.
    
    \item \textbf{Visualisatie}: de structuur wordt omgezet naar Mermaid.js-syntax, die vervolgens wordt omgezet naar een flowchart.
\end{enumerate}

\newpage

\subsection{Stap 1: Parsing}

In deze eerste stap wordt de HLASM-code verwerkt tot een lijst van instructies.
Lege lijnen en commentaar (lijnen die starten met een *) worden genegeerd.

Voor elke lijn worden drie onderdelen bepaald:
\begin{itemize}
    \item het label (indien aanwezig),
    \item de mnemonic (de instructienaam),
    \item het eerste operand (argument).
\end{itemize}

De parser maakt gebruik van de vaste kolomstructuur door specifieke delen van de lijn uit te lezen op vaste posities.

Na deze stap wordt elke lijn voorgesteld als een object met duidelijke velden.  
Dit maakt het eenvoudiger om de volgende stap uit te voeren.

\begin{lstlisting}[language=JavaScript, caption={De HLASM parser.}]
function parseInstructions(hlasmCode) {
    const instructions = [];
    
    hlasmCode.split('\n').forEach((line, index) => {
        if (!line || line.trim() === "" || line.trim().startsWith("*"))
        return;
        
        const padded   = line.padEnd(80, " ");
        const label    = padded.substring(0, 8).trim();
        const mnemonic = padded.substring(9, 14).trim().toUpperCase();
        
        if (!mnemonic) return;
        
        const operandField = padded.substring(15, 71).trim();
        const operandMatch = operandField.match(/^([^\s']+|'[^']*')+/);
        const fullOperand  = operandMatch ? operandMatch[0] : "";
        const firstOperand = fullOperand.split(',')[0].trim();
        
        instructions.push({
            lineNumber: index + 1,
            raw: `${label.padEnd(8)} ${mnemonic.padEnd(5)} ${fullOperand}`,
            label,
            mnemonic,
            operand: firstOperand
        });
    });
    
    return instructions;
}
\end{lstlisting}

\subsection{Stap 2: Control-flow analyse}

In deze stap wordt gekeken hoe het programma verloopt tijdens uitvoering.

De instructies worden gegroepeerd in blokken.  
Een blok is een reeks instructies die na elkaar uitgevoerd worden, zonder onderbreking.

Een nieuw blok start in de volgende situaties:
\begin{itemize}
    \item wanneer een instructie een label bevat,
    \item wanneer er een sprong (branch) gebeurt,
    \item wanneer het programma eindigt (END).
\end{itemize}

Voor elk blok wordt ook bepaald hoe het blok eindigt.  
Dit kan bijvoorbeeld een sprong zijn naar een ander label, of een gewone overgang naar het volgende blok.

Op deze manier ontstaat een structuur die beschrijft hoe de uitvoering van het programma kan verlopen.

\begin{lstlisting}[language=JavaScript, caption={De code die instructies in blokken verdeelt.}]
function buildBlocks(instructions) {
    let blockId = 0;
    const blocks = [];
    let currentBlock = createBlock(0);
    
    instructions.forEach((instr, i) => {
        if (instr.label && currentBlock.lines.length > 0) {
            blocks.push(currentBlock);
            currentBlock = createBlock(++blockId);
        }
        
        currentBlock.lines.push(instr);
        
        if (instr.label) {
            currentBlock.labels.push(instr.label);
        }
        
        const isBranch = isBranchInstruction(instr.mnemonic);
        const isEnd = instr.mnemonic === "END";
        
        if (isBranch || isEnd) {
            currentBlock.exit = instr;
            blocks.push(currentBlock);
            currentBlock = createBlock(++blockId);
        } else if (instructions[i + 1] && instructions[i + 1].label) {
            currentBlock.exit = { mnemonic: "SEQUENTIAL", operand: "" };
            blocks.push(currentBlock);
            currentBlock = createBlock(++blockId);
        }
    });
    
    if (currentBlock.lines.length > 0 && !blocks.includes(currentBlock)) {
        blocks.push(currentBlock);
    }
    
    return blocks;
}
\end{lstlisting}

\subsection{Stap 3: Visualisatie}

In de laatste stap wordt de gevonden structuur omgezet naar een visueel diagram.

Eerst worden alle blokken voorgesteld als knopen (nodes).  
De inhoud van elk blok wordt weergegeven als tekst in de knoop.

Daarna worden verbindingen (pijlen) toegevoegd tussen de blokken:
\begin{itemize}
    \item een gewone pijl voor een normale overgang,
    \item een pijl met label voor een conditionele sprong,
    \item een directe pijl voor een onvoorwaardelijke sprong.
\end{itemize}

Deze informatie wordt omgezet naar Mermaid.js-syntaxis.  
Mermaid.js zorgt er vervolgens voor dat deze tekst automatisch wordt weergegeven als een flowchart.

\begin{lstlisting}[language=JavaScript, caption={De conversie naar mermaid.js.}]
function blocksToMermaid(blocks) {
    let mermaid = "%%{init: {'themeVariables': {'fontFamily': 'monospace'}}}%%\nflowchart TD\n";
    
    blocks.forEach(block => {
        const text = block.lines
        .map(l => l.raw.replace(/\s+/, ' ').trim()).join("<br/>");
        mermaid += `    ${block.id}["${text}"]\n`;
    });
    
    blocks.forEach((block, i) => {
        const next = blocks[i + 1];
        const exit = block.exit;
        
        // we have reached the last block
        if (!exit) return;
        
        const m = exit.mnemonic;
        const targetLabel = exit.operand;
        
        if (isConditionalBranch(m)) {
            const targetBlock = findBlockByLabel(blocks, targetLabel);
            if (targetBlock) mermaid += `    ${block.id} -->|${m}| ${targetBlock.id}\n`;
            if (next) mermaid += `    ${block.id} --> ${next.id}\n`;
        } else if (isUnconditionalBranch(m)) {
            const targetBlock = findBlockByLabel(blocks, targetLabel);
            if (targetBlock) mermaid += `    ${block.id} --> ${targetBlock.id}\n`;
        } else if (m === "SEQUENTIAL" && next) {
            mermaid += `    ${block.id} --> ${next.id}\n`;
        }
    });
    
    return mermaid;
}
\end{lstlisting}

\subsection{Gebruikersinterface}

De gebruikersinterface van de applicatie is eenvoudig opgebouwd.

Er zijn twee onderdelen:
\begin{itemize}
    \item een invoerveld waarin de gebruiker HLASM-code kan invoeren,
    \item een visualisatiepaneel waarin het diagram wordt weergegeven.
\end{itemize}

Wanneer de gebruiker code invoert, wordt het diagram automatisch bijgewerkt.  
De gebruiker ziet dus onmiddellijk het resultaat van de analyse.

\begin{figure}[h!]
    \centering
    \includegraphics[width=0.9\textwidth]{gebruikersinterface.png}
    \caption{Invoer van HLASM-code in de webapplicatie.}
\end{figure}

\newpage

\subsection{Gebruik van de applicatie}

Het gebruik van de tool verloopt in enkele eenvoudige stappen:

\begin{enumerate}
    \item De gebruiker plakt of typt HLASM-broncode in de editor.
    \item De applicatie verwerkt de code automatisch.
    \item Op basis van de analyse wordt een Mermaid.js-diagram gegenereerd.
    \item Het diagram wordt onmiddellijk visueel weergegeven.
\end{enumerate}

De gebruiker hoeft geen kennis te hebben van Mermaid.js. De visualisatie wordt volledig automatisch gegenereerd en helpt bij het begrijpen van de structuur van het programma.

\section{Beperkingen van de implementatie}

Omdat de implementatie gebruik maakt van een vereenvoudige aanpak zijn er enkele beperkingen.

Ten eerste wordt geen volledige syntactische of semantische analyse uitgevoerd. De parser gaat ervan uit dat de invoer de standaard kolomstructuur volgt. Als dit niet het geval is, kan de interpretatie fout zijn.

Daarnaast worden macro-instructies niet uitgewerkt. Macros kunnen meerdere instructies bevatten, waardoor een deel van de control-flow niet zichtbaar is in het diagram.

Verder worden enkel directe sprongen naar labels ondersteund.  
Indirecte sprongen, bijvoorbeeld via registers (zoals \texttt{BR R14}), worden niet geanalyseerd.

Tot slot wordt de condition code niet inhoudelijk geïnterpreteerd.  

Deze beperkingen zijn aanvaardbaar binnen deze proof-of-concept.  
De focus ligt op het zichtbaar maken van de globale structuur van het programma.

\newpage

\section{Evaluatie}

De evaluatie onderzoekt of de tool helpt bij het begrijpen van HLASM-programma’s.

Alle deelnemers hadden geen voorkennis van HLASM.  
Ze kregen vooraf een korte uitleg over control-flow, met nadruk op labels en sprongen.

De evaluatie werd uitgevoerd met vier deelnemers.  
Elke deelnemer kreeg twee codefragmenten en moest:
\begin{itemize}
    \item het startpunt identificeren,
    \item vertakkingen aanduiden,
    \item de globale structuur beschrijven.
\end{itemize}

Deze opdrachten werden uitgevoerd:
\begin{itemize}
    \item zonder hulpmiddel,
    \item met gebruik van de tool.
\end{itemize}

\subsection{Resultaten}

Uit de evaluatie blijkt dat de tool vooral helpt om snel een overzicht te krijgen van de globale structuur van een programma.

De deelnemers konden met de visualisatie sneller zien hoe de verschillende delen van het programma met elkaar verbonden zijn en waar vertakkingen plaatsvinden. Dit maakte het eenvoudiger om de algemene flow te begrijpen.

Tegelijk werd duidelijk dat dit overzicht niet voldoende is om de code volledig te begrijpen.  
De deelnemers gaven aan dat het nog steeds nodig is om de volledige code lijn per lijn te bekijken om de werking correct te interpreteren.

De tool werd daarom niet gezien als een vervanging van manuele analyse, maar als een hulpmiddel dat het proces ondersteunt.

Daarnaast werd opgemerkt dat de meerwaarde van de tool vooral zichtbaar is bij grotere programma’s.  
Bij kleinere codefragmenten is de structuur meestal eenvoudig genoeg om manueel te volgen, waardoor de visualisatie weinig extra toevoegt.
